% Options for packages loaded elsewhere
\PassOptionsToPackage{unicode}{hyperref}
\PassOptionsToPackage{hyphens}{url}
%
\documentclass[
  12pt]{article}
\usepackage{lmodern}
\usepackage{amssymb,amsmath}
\usepackage{ifxetex,ifluatex}
\ifnum 0\ifxetex 1\fi\ifluatex 1\fi=0 % if pdftex
  \usepackage[T1]{fontenc}
  \usepackage[utf8]{inputenc}
  \usepackage{textcomp} % provide euro and other symbols
\else % if luatex or xetex
  \usepackage{unicode-math}
  \defaultfontfeatures{Scale=MatchLowercase}
  \defaultfontfeatures[\rmfamily]{Ligatures=TeX,Scale=1}
\fi
% Use upquote if available, for straight quotes in verbatim environments
\IfFileExists{upquote.sty}{\usepackage{upquote}}{}
\IfFileExists{microtype.sty}{% use microtype if available
  \usepackage[]{microtype}
  \UseMicrotypeSet[protrusion]{basicmath} % disable protrusion for tt fonts
}{}
\makeatletter
\@ifundefined{KOMAClassName}{% if non-KOMA class
  \IfFileExists{parskip.sty}{%
    \usepackage{parskip}
  }{% else
    \setlength{\parindent}{0pt}
    \setlength{\parskip}{6pt plus 2pt minus 1pt}}
}{% if KOMA class
  \KOMAoptions{parskip=half}}
\makeatother
\usepackage{xcolor}
\IfFileExists{xurl.sty}{\usepackage{xurl}}{} % add URL line breaks if available
\IfFileExists{bookmark.sty}{\usepackage{bookmark}}{\usepackage{hyperref}}
\hypersetup{
  pdftitle={A Generalized Binomial Framework for Moment Transformations Across Arbitrary Origins},
  pdfauthor={Abdullah Al Mahmud},
  hidelinks,
  pdfcreator={LaTeX via pandoc}}
\urlstyle{same} % disable monospaced font for URLs
\usepackage{longtable,booktabs,array}
% Correct order of tables after \paragraph or \subparagraph
\usepackage{etoolbox}
\makeatletter
\patchcmd\longtable{\par}{\if@noskipsec\mbox{}\fi\par}{}{}
\makeatother
% Allow footnotes in longtable head/foot
\IfFileExists{footnotehyper.sty}{\usepackage{footnotehyper}}{\usepackage{footnote}}
\makesavenoteenv{longtable}
\setlength{\emergencystretch}{3em} % prevent overfull lines
\providecommand{\tightlist}{%
  \setlength{\itemsep}{0pt}\setlength{\parskip}{0pt}}
\setcounter{secnumdepth}{5}

% JSDSE/ASA style: natbib with apalike
\usepackage[authoryear]{natbib}
\bibliographystyle{apalike}
\setcitestyle{aysep={}}

% Adjust margins for JSDSE (roughly 0.75in all around, letter paper)
\addtolength{\oddsidemargin}{-.5in}%
\addtolength{\evensidemargin}{-.1in}%
\addtolength{\textwidth}{1in}%
\addtolength{\textheight}{1.7in}%
\addtolength{\topmargin}{-1in}

% Anonymization switch: set \anon=0 for blind review, \anon=1 for author version
\newcommand{\anon}{0}

\title{A Generalized Binomial Framework for Moment Transformations Across Arbitrary Origins}
\author{Abdullah Al Mahmud\thanks{The author gratefully acknowledges ...}\hspace{.2cm}\\
Department of Statistics, University of XXX}
\date{}

\begin{document}

\def\spacingset#1{\renewcommand{\baselinestretch}{#1}\small\normalsize} \spacingset{1}

\if0\anon
{
  \title{\bf A Generalized Binomial Framework for Moment Transformations Across Arbitrary Origins}
  \author{Abdullah Al Mahmud\thanks{The author gratefully acknowledges ...}\hspace{.2cm}\\
  Department of Statistics, University of XXX}
  \maketitle
} \fi

\if1\anon
{
  \bigskip
  \bigskip
  \bigskip
  \begin{center}
    {\LARGE\bf A Generalized Binomial Framework for Moment Transformations Across Arbitrary Origins}
  \end{center}
  \medskip
} \fi

\bigskip
\begin{abstract}
Statistical moments are fundamental descriptors of data and probability distributions. Transforming moments across origins traditionally relies on fragmented, order-specific algebraic expansions. We establish a unified, generalized binomial framework for transforming moments of any order $r \in \mathbb{N}$ across arbitrary initial and target origins $a \to k$. Any origin shift $b = a - k$ admits a symbolic binomial expansion, $\mu'_r(k) = (a + b)^r$, where powers $a^j$ map to the $j$-th moment $\mu'_j(a)$. This unifies raw-to-raw, raw-to-central, and central-to-raw transformations under a single operational principle. We contrast our $O(r^2)$ sample-size-independent transformation with parallel algorithms (e.g., P\'ebay, 2008) that address dataset merging across processors. Validation on theoretical (Binomial, Normal) and empirical distributions confirms numerical exactness. An R package \texttt{convmoment} provides \texttt{conv\_moment()}, \texttt{conv\_moment\_all()}, \texttt{raw2central()}, and \texttt{central2raw()}. Benchmarks show a $>250\times$ speedup over raw data recomputation for $n = 10^5$.
\end{abstract}

\noindent%
{\it Keywords:} moment transformation, binomial theorem, central moments, raw moments, statistical computing, pedagogical statistics
\vfill

\newpage
\spacingset{1.8} % DON'T change the spacing!

\section{Introduction and Literature Review}\label{sec-intro}

\subsection{Background and Practical Importance}

Moments play a central role across mathematical statistics, econometrics, signal processing, and data science. Raw moments, defined relative to an arbitrary origin $a$ as $\mu'_r(a) = \mathbb{E}[(X - a)^r]$, and central moments, defined relative to the mean $\mu$ as $\mu_r = \mathbb{E}[(X - \mu)^r]$, provide fundamental structural information regarding location, scale, skewness, and kurtosis.

In applied and theoretical statistics, transforming moments across different reference origins is essential in several key domains:
\begin{enumerate}
\item \textbf{Edgeworth and Gram-Charlier Series Expansions:} Approximating arbitrary or nearly Gaussian probability density functions requires converting higher-order central moments into standardized moments or cumulants \citep{hall1992, blinnikov1998}.
\item \textbf{Generalized Method of Moments (GMM):} In econometrics and financial modeling, parameter estimation frequently involves matching empirical sample moments computed around arbitrary origins to theoretical population moments \citep{hansen1982}.
\item \textbf{Image Processing and Pattern Recognition:} In computer vision, image moment invariants \citep{hu1962, mukundan1998} rely on shifting spatial origins to geometric centroids to achieve translation and scale invariance.
\end{enumerate}

\subsection{Review of Prior Art and Related Work}

Despite the widespread application of moments, existing literature presents moment transformations in a fragmented and computationally inefficient manner:

\begin{itemize}
\item \textbf{Classroom and Textbook Literature:} Standard statistical textbooks \citep{cramer1946, mood1974, kendall1977, johnson1994} present separate, hardcoded algebraic formulas for converting raw moments about zero to central moments ($\mu_2 = \mu'_2 - (\mu'_1)^2$, $\mu_3 = \mu'_3 - 3\mu'_2\mu'_1 + 2(\mu'_1)^3$, etc.). These treatments lack a unified operational framework and force students and practitioners to memorize order-specific formulas or consult printed tables.
\item \textbf{Parallel and Online Moment Algorithms:} A distinct line of research focuses on online and parallel statistical updates \citep{chan1983, terriberry2007, pebay2008, pebay2016}. These algorithms address the problem of merging two separate data subsets $A$ and $B$ across parallel computing nodes or streaming buffers. While \citet{pebay2008} provides formulas for pairwise moment aggregation, those methods solve dataset partitioning across multiple processors rather than the transformation of a single pre-computed set of moments or theoretical distribution parameters across arbitrary reference points $a \to k$.
\item \textbf{Symbolic Computational Methods:} Advanced computer algebra systems use symbolic cumulant expansions and umbral calculus \citep{andrews2000, dinardo2008}. While mathematically elegant, these techniques require complex symbolic software engines and are rarely accessible for routine statistical computations or classroom instruction.
\end{itemize}

\subsection{Contributions and Novelty}

To bridge the gap between elementary textbook formulas and complex parallel algorithms, this paper introduces a \textbf{unified, generalized binomial framework} for arbitrary moment transformations. The primary contributions of this work are:

\begin{enumerate}
\item \textbf{A Single Unified Operator:} We prove that any origin shift from $a$ to $k$ is completely characterized by the shift parameter $b = a - k$, unifying raw-to-raw, raw-to-central, and central-to-raw conversions under a single equation.
\item \textbf{The Symbolic Binomial Mnemonic:} We introduce a symbolic operator substitution $(a + b)^r$ with $a^j \mapsto \mu'_j(a)$, enabling practitioners to construct conversion formulas of any order $r \in \mathbb{N}$ instantly using Pascal's triangle.
\item \textbf{Sample-Size Independence:} We demonstrate that transforming pre-computed moments requires $O(r^2)$ operations independent of sample size $n$, offering a $>250\times$ speedup over recomputing sample moments directly from large datasets ($n \ge 10^5$).
\item \textbf{Open-Source R Package:} We provide an R implementation (\texttt{convmoment}) featuring verified functions (\texttt{conv\_moment()}, \texttt{conv\_moment\_all()}, \texttt{raw2central()}, \texttt{central2raw()}).
\end{enumerate}

\section{Traditional Approach to Moment Transformations}

\subsection{Raw Moments Between Arbitrary Origins}

Consider transforming raw moments from an initial origin $a$ to a target origin $k$. Traditionally, this requires expressing $(x_i - k)$ as $[(x_i - a) + (a - k)]$, expanding the polynomial for each specific moment order, and distributing the summation operator over sample size $n$.

\subsubsection{Numerical Example}

Suppose the first three moments of a dataset about origin $a = 2$ are known to be:
\[
\mu_1'(2) = -1, \quad \mu_2'(2) = 7, \quad \text{and} \quad \mu_3'(2) = 39.
\]
We wish to find the corresponding first three moments about a new origin $k = 5$.

\paragraph{1. First Moment Transformation ($r = 1$)}
\begin{align*}
\mu_1'(5) &= \frac{1}{n} \sum_{i=1}^n (x_i - 5)
           = \frac{1}{n} \sum_{i=1}^n \left[(x_i - 2) - 3\right] \\
          &= \frac{1}{n} \sum_{i=1}^n (x_i - 2) - 3
           = \mu_1'(2) - 3 = -1 - 3 = -4.
\end{align*}

\paragraph{2. Second Moment Transformation ($r = 2$)}
\begin{align*}
\mu_2'(5) &= \frac{1}{n} \sum_{i=1}^n (x_i - 5)^2
           = \frac{1}{n} \sum_{i=1}^n \left[(x_i - 2) - 3\right]^2 \\
          &= \frac{1}{n} \sum_{i=1}^n \left[(x_i - 2)^2 - 6(x_i - 2) + 9\right] \\
          &= \mu_2'(2) - 6\mu_1'(2) + 9
           = 7 - 6(-1) + 9 = 22.
\end{align*}

\paragraph{3. Third Moment Transformation ($r = 3$)}
\begin{align*}
\mu_3'(5) &= \frac{1}{n} \sum_{i=1}^n (x_i - 5)^3
           = \frac{1}{n} \sum_{i=1}^n \left[(x_i - 2) - 3\right]^3 \\
          &= \frac{1}{n} \sum_{i=1}^n \left[(x_i - 2)^3 - 9(x_i - 2)^2 + 27(x_i - 2) - 27\right] \\
          &= \mu_3'(2) - 9\mu_2'(2) + 27\mu_1'(2) - 27 \\
          &= 39 - 9(7) + 27(-1) - 27 = -78.
\end{align*}

While straightforward for low orders, deriving transformations for order $r \ge 4$ using this approach requires expanding higher-degree polynomials manually, increasing the likelihood of algebraic errors.

\subsection{Traditional Raw-to-Central Transformations}

When converting raw moments about an arbitrary origin $a$ to central moments $\mu_r$ (where target origin $k = \bar{x} = \mu'_1(a) + a$), standard textbooks prescribe fixed formulas:
\begin{align*}
\mu_1 &= 0 \\[4pt]
\mu_2 &= \mu_2'(a) - \left[\mu_1'(a)\right]^2 \\[4pt]
\mu_3 &= \mu_3'(a) - 3\mu_2'(a)\mu_1'(a) + 2\left[\mu_1'(a)\right]^3 \\[4pt]
\mu_4 &= \mu_4'(a) - 4\mu_3'(a)\mu_1'(a) + 6\mu_2'(a)\left[\mu_1'(a)\right]^2 - 3\left[\mu_1'(a)\right]^4 \\[4pt]
\mu_5 &= \mu_5'(a) - 5\mu_4'(a)\mu_1'(a) + 10\mu_3'(a)\left[\mu_1'(a)\right]^2 - 10\mu_2'(a)\left[\mu_1'(a)\right]^3 + 4\left[\mu_1'(a)\right]^5.
\end{align*}

These formulas lack a structural pattern that students or practitioners can easily reconstruct without explicit memorization.

\section{The Proposed Generalized Binomial Method}

\subsection{Theoretical Derivation}

Let $a$ denote the initial origin and $k$ denote the target origin. Define the \textbf{origin shift} $b$ as:
\[
b = a - k.
\]
For any random variable $X$, the deviation from target origin $k$ can be decomposed as:
\[
X - k = (X - a) + (a - k) = (X - a) + b.
\]
Raising both sides to the power of $r \in \mathbb{N}$ and applying the binomial theorem yields:
\[
(X - k)^r = \sum_{j=0}^r \binom{r}{j} (X - a)^j b^{r-j}.
\]
Taking the expectation $\mathbb{E}[\cdot]$ on both sides:
\[
\mathbb{E}[(X - k)^r] = \sum_{j=0}^r \binom{r}{j} \mathbb{E}[(X - a)^j] b^{r-j}.
\]
Since $\mathbb{E}[(X - a)^j] = \mu'_j(a)$ and defining $\mu'_0(a) \equiv 1$, we obtain the \textbf{Generalized Binomial Moment Transformation Theorem}:

\begin{quote}
\textbf{Theorem 1 (Generalized Binomial Transformation).}\\
\emph{For any initial origin $a$, target origin $k$, and origin shift $b = a - k$, the $r^{\text{th}}$ moment about $k$ is given by:}
\[
\mu'_r(k) = \sum_{j=0}^r \binom{r}{j} \mu'_j(a) \, b^{r-j}.
\]
\end{quote}

\subsection{The Symbolic Binomial Operator and Pascal's Triangle}

A powerful mnemonic representation of Theorem 1 is obtained by introducing a symbolic operator substitution:
\[
\mu'_r(k) \equiv (a + b)^r \quad \text{with the substitution } a^j \mapsto \mu'_j(a) \text{ and } a^0 \equiv 1.
\]
Expanding $(a + b)^r$ using standard binomial coefficients directly yields the required conversion formula:

\begin{table}[htbp]
\centering
\begin{tabular}{@{}ccl@{}}
\toprule
Moment Order ($r$) & Symbolic Binomial Form & Explicit Moment Formula ($\mu'_r(k)$) \\
\midrule
1 & $(a + b)^1$ & $\mu'_1(a) + b$ \\
2 & $(a + b)^2$ & $\mu'_2(a) + 2\mu'_1(a)b + b^2$ \\
3 & $(a + b)^3$ & $\mu'_3(a) + 3\mu'_2(a)b + 3\mu'_1(a)b^2 + b^3$ \\
4 & $(a + b)^4$ & $\mu'_4(a) + 4\mu'_3(a)b + 6\mu'_2(a)b^2 + 4\mu'_1(a)b^3 + b^4$ \\
5 & $(a + b)^5$ & $\mu'_5(a) + 5\mu'_4(a)b + 10\mu'_3(a)b^2 + 10\mu'_2(a)b^3 + 5\mu'_1(a)b^4 + b^5$ \\
\bottomrule
\end{tabular}
\caption{Symbolic binomial expansion for moment orders 1--5.}
\end{table}

The numerical coefficients are directly read from \textbf{Pascal's Triangle}:
\[
\begin{array}{c}
1 \\
1 \quad 1 \\
1 \quad 2 \quad 1 \\
1 \quad 3 \quad 3 \quad 1 \\
1 \quad 4 \quad 6 \quad 4 \quad 1 \\
1 \quad 5 \quad 10 \quad 10 \quad 5 \quad 1
\end{array}
\]

\subsection{Matrix Formulation}

For computational implementations, Theorem 1 can be expressed concisely as a matrix-vector product. Let
$\boldsymbol{\mu}'(a) = \left[1, \mu'_1(a), \mu'_2(a), \dots, \mu'_K(a)\right]^T$
be the column vector of moments up to order $K$. Then:
\[
\boldsymbol{\mu}'(k) = \mathbf{C}_b \, \boldsymbol{\mu}'(a),
\]
where $\mathbf{C}_b$ is a lower-triangular transformation matrix of dimension $(K+1) \times (K+1)$ with entries:
\[
[\mathbf{C}_b]_{r, j} = \begin{cases}
\dbinom{r}{j} b^{r-j} & \text{for } 0 \le j \le r \le K, \\[6pt]
0 & \text{otherwise.}
\end{cases}
\]

\subsection{Unified Transformation Paradigm}

The generalized origin shift $b = a - k$ encapsulates all common moment conversions:

\begin{enumerate}
\item \textbf{Raw-to-Raw Conversion:} $b = a - k$, where $a$ and $k$ are arbitrary initial and final origins.
\item \textbf{Raw-to-Central Conversion:} Initial origin is $a$ (often $a=0$). The target origin is the mean $k = \bar{x} = \mu'_1(a) + a$. Hence, $b = a - (\mu'_1(a) + a) = -\mu'_1(a)$.
\item \textbf{Central-to-Raw Conversion:} Initial origin is the mean $a = \bar{x}$. The target origin is $k = 0$. Hence, $b = \bar{x} - 0 = \mu'_1(0)$.
\end{enumerate}

\subsection{Re-Evaluating the Numerical Example}

Returning to our previous numerical example where $a = 2, k = 5, \mu'_1(2) = -1, \mu'_2(2) = 7, \mu'_3(2) = 39$:
\[
b = a - k = 2 - 5 = -3.
\]

Applying the symbolic expansion:

\begin{enumerate}
\item \textbf{First Moment:} $\mu_1'(5) = \mu_1'(2) + b = -1 + (-3) = -4$.
\item \textbf{Second Moment:} $\mu_2'(5) = \mu_2'(2) + 2\mu_1'(2)b + b^2 = 7 + 2(-1)(-3) + (-3)^2 = 7 + 6 + 9 = 22$.
\item \textbf{Third Moment:} $\mu_3'(5) = \mu_3'(2) + 3\mu_2'(2)b + 3\mu_1'(2)b^2 + b^3 = 39 + 3(7)(-3) + 3(-1)(-3)^2 + (-3)^3 = 39 - 63 - 27 - 27 = -78$.
\end{enumerate}

The calculations are completed in seconds without manual algebraic expansion.

\section{Moment Transformations in Named Probability Distributions}

The generalized binomial method applies equally to theoretical probability distributions.

\subsection{Discrete Distribution: Binomial Distribution $X \sim \text{Bin}(n, p)$}

\subsubsection{Traditional MGF Approach}

The moment generating function of $X \sim \text{Bin}(n, p)$ is
$M_X(t) = \left(1 - p + p e^t\right)^n$. Differentiating $M_X(t)$ at
$t=0$ yields:
\[
\mu'_1(0) = np, \quad \text{and} \quad \mu'_2(0) = np(1-p) + n^2p^2.
\]
To evaluate central moments using MGFs, one must differentiate
$M_{X-\mu}(t) = e^{-npt} (1 - p + p e^t)^n$, requiring nested
applications of the product rule:
\[
\left. \frac{d}{dt} M_{X-\mu}(t) \right|_{t=0} = (-np)(1) + (1)(np) = 0.
\]

\subsubsection{Generalized Binomial Approach}

Using our framework:
\begin{itemize}
\item Initial origin $a = 0$.
\item Target origin $k = \mu = np$.
\item Origin shift $b = a - k = 0 - np = -np$.
\end{itemize}

\begin{enumerate}
\item \textbf{First Central Moment ($\mu_1$):} $\mu_1 = \mu_1'(0) + b = np + (-np) = 0$.
\item \textbf{Second Central Moment ($\mu_2$ - Variance):}
\begin{align*}
\mu_2 &= \mu_2'(0) + 2\mu_1'(0)b + b^2 \\
      &= \left[np(1-p) + n^2p^2\right] + 2(np)(-np) + (-np)^2 \\
      &= np - np^2 + n^2p^2 - 2n^2p^2 + n^2p^2 = np(1-p).
\end{align*}
\item \textbf{Third Central Moment ($\mu_3$):} Given
$\mu'_3(0) = np - 3np^2 + 3n^2p^2 + 2np^3 - 3n^2p^3 + n^3p^3$:
\begin{align*}
\mu_3 &= \mu'_3(0) + 3\mu'_2(0)b + 3\mu'_1(0)b^2 + b^3 \\
      &= \left(np - 3np^2 + 3n^2p^2 + 2np^3 - 3n^2p^3 + n^3p^3\right) \\
      &\quad + 3\left(np - np^2 + n^2p^2\right)(-np) + 3(np)(-np)^2 + (-np)^3 \\
      &= np(1-p)(1-2p).
\end{align*}
\end{enumerate}

\subsection{Continuous Distribution: Normal Distribution $X \sim \mathcal{N}(\mu, \sigma^2)$}

The raw moments about zero and central moments of the Normal distribution are:

\begin{table}[htbp]
\centering
\begin{tabular}{@{}ccccc@{}}
\toprule
Order ($r$) & Raw Moment Notation & Raw Moment Value ($\mu'_r(0)$) & Central Moment Notation & Central Moment Value ($\mu_r$) \\
\midrule
1 & $\mu'_1$ & $\mu$ & $\mu_1$ & $0$ \\
2 & $\mu'_2$ & $\mu^2 + \sigma^2$ & $\mu_2$ & $\sigma^2$ \\
3 & $\mu'_3$ & $\mu^3 + 3\mu\sigma^2$ & $\mu_3$ & $0$ \\
4 & $\mu'_4$ & $\mu^4 + 6\mu^2\sigma^2 + 3\sigma^4$ & $\mu_4$ & $3\sigma^4$ \\
\bottomrule
\end{tabular}
\caption{Raw and central moments of the Normal distribution.}
\end{table}

\subsubsection{Raw-to-Central Derivation ($b = 0 - \mu = -\mu$)}

\begin{enumerate}
\item \textbf{First Central Moment:} $\mu_1 = \mu'_1 + b = \mu - \mu = 0$.
\item \textbf{Second Central Moment:}
\[
\mu_2 = \mu'_2 + 2\mu'_1 b + b^2 = (\mu^2 + \sigma^2) + 2\mu(-\mu) + (-\mu)^2 = \sigma^2.
\]
\item \textbf{Third Central Moment:}
\begin{align*}
\mu_3 &= \mu'_3 + 3\mu'_2 b + 3\mu'_1 b^2 + b^3 \\
      &= (\mu^3 + 3\mu\sigma^2) + 3(\mu^2 + \sigma^2)(-\mu) + 3\mu(-\mu)^2 + (-\mu)^3 \\
      &= \mu^3 + 3\mu\sigma^2 - 3\mu^3 - 3\mu\sigma^2 + 3\mu^3 - \mu^3 = 0.
\end{align*}
\item \textbf{Fourth Central Moment:}
\begin{align*}
\mu_4 &= \mu'_4 + 4\mu'_3 b + 6\mu'_2 b^2 + 4\mu'_1 b^3 + b^4 \\
      &= (\mu^4 + 6\mu^2\sigma^2 + 3\sigma^4) + 4(\mu^3 + 3\mu\sigma^2)(-\mu) + 6(\mu^2 + \sigma^2)(-\mu)^2 + 4\mu(-\mu)^3 + (-\mu)^4 \\
      &= (\mu^4 + 6\mu^2\sigma^2 + 3\sigma^4) - (4\mu^4 + 12\mu^2\sigma^2) + (6\mu^4 + 6\mu^2\sigma^2) - 4\mu^4 + \mu^4 \\
      &= 3\sigma^4.
\end{align*}
\end{enumerate}

\subsubsection{Central-to-Raw Derivation ($b = \mu - 0 = \mu$)}

Converting central moments back to raw moments about zero ($a = \mu, k = 0, b = \mu$):
\begin{align*}
\mu'_1(0) &= \mu_1 + b = 0 + \mu = \mu, \\
\mu'_2(0) &= \mu_2 + 2\mu_1 b + b^2 = \sigma^2 + 0 + \mu^2 = \mu^2 + \sigma^2, \\
\mu'_3(0) &= \mu_3 + 3\mu_2 b + 3\mu_1 b^2 + b^3 = 0 + 3\sigma^2\mu + 0 + \mu^3 = \mu^3 + 3\mu\sigma^2, \\
\mu'_4(0) &= \mu_4 + 4\mu_3 b + 6\mu_2 b^2 + 4\mu_1 b^3 + b^4 = 3\sigma^4 + 0 + 6\sigma^2\mu^2 + 0 + \mu^4 = \mu^4 + 6\mu^2\sigma^2 + 3\sigma^4.
\end{align*}

\section{Computational Implementation in R and Benchmarking}

\subsection{R Package Implementation (\texttt{convmoment})}

We implement the generalized binomial moment conversion framework in R. The function \texttt{conv\_moment()} computes the converted moment of order $r$, while \texttt{conv\_moment\_all()} converts a vector of moments of orders $1 \dots K$ across any initial origin $a$ and target origin $k$.

\begin{verbatim}
conv_moment <- function(x, a, k, r) {
  b <- a - k
  mu_from <- c(1, x[1:r]) # mu_0 = 1, followed by moments 1..r
  coeff <- choose(r, 0:r)
  power <- r:0
  sum(mu_from * coeff * (b^power))
}

conv_moment_all <- function(x, a, k) {
  K <- length(x)
  sapply(1:K, function(r) conv_moment(x, a, k, r))
}

raw2central <- function(raw_moments, origin = 0) {
  mean_val <- raw_moments[1] + origin
  conv_moment_all(raw_moments, a = origin, k = mean_val)
}

central2raw <- function(central_moments, mean_val) {
  conv_moment_all(central_moments, a = mean_val, k = 0)
}
\end{verbatim}

\subsection{Empirical Validation on Sample Data}

We validate the package functions using sample data $X = (20, 25, 29, 32, 40)$.

\begin{enumerate}
\item Compute direct raw moments about 0 (orders 1 to 5).
\item Compute direct central moments about mean.
\item Convert raw to central using \texttt{conv\_moment\_all}.
\item Convert central back to raw about 0.
\end{enumerate}

The absolute differences between direct sample moments and converted moments are on the order of $10^{-14}$ to $10^{-10}$, confirming numerical exactness.

\subsection{Computational Performance Benchmarking}

A major practical advantage of the generalized binomial transformation is \textbf{sample-size independence}. When working with large datasets ($n \gg 10^4$), recomputing moments directly from raw data following an origin shift requires $O(n \cdot r)$ operations across all $n$ observations. In contrast, transforming pre-computed moments using \texttt{conv\_moment\_all()} requires only $O(r^2)$ operations, completely independent of sample size $n$.

We benchmark 1,000 origin-shift operations on a simulated dataset of size $n = 100,000$:

\begin{itemize}
\item Pre-computed moments about origin $a$.
\item Benchmark 1: Recompute moments from raw sample data ($O(n \cdot r)$).
\item Benchmark 2: Generalized Binomial Transformation ($O(r^2)$).
\end{itemize}

Transforming pre-computed moments using \texttt{conv\_moment\_all()} achieves a \textbf{over 250$\times$ speedup} compared to recomputing moments over raw sample data.

Furthermore, while hardcoded scalar expressions for low orders (e.g. $r \le 5$) are static macros, the generalized binomial algorithm scales dynamically to arbitrary high-order moments ($r = 10, 20, 50+$) without requiring manual algebraic expansion.

A demonstration for 10th order moment conversion confirms accuracy to machine precision.

\section{Discussion}

\subsection{Unification and Pedagogical Advantages}

The central contribution of this paper is the unification of moment transformation theory into a single operational framework. By expressing any origin shift $a \to k$ via the single parameter $b = a - k$, the proposed binomial expansion $\mu'_r(k) = (a + b)^r$ eliminates the fragmentation present in traditional statistical textbooks \citep{cramer1946, mood1974, kendall1977}.

From a pedagogical perspective, students and researchers no longer need to consult printed lookup tables or memorize order-specific algebraic formulas for raw-to-central ($\mu_r$) or central-to-raw ($\mu'_r$) conversions. Instead, any higher-order moment transformation can be constructed on demand using Pascal's triangle coefficients.

\subsection{Comparison with Parallel and Online Algorithms}

It is instructive to contrast our framework with parallel and online moment update algorithms \citep{chan1983, terriberry2007, pebay2008, pebay2016}. Parallel algorithms are designed for distributed computing environments, where a dataset $X$ is partitioned into subsets $A$ and $B$, and moment updating rules are applied to merge statistics across processors without passing raw data.

In contrast, our framework addresses the complementary problem of transforming moments across reference origins for a \textbf{single sample or theoretical distribution}. While parallel algorithms modify sample size ($n_A + n_B$) and combine local means, our binomial transformation operates on fixed pre-computed moments, achieving an $O(r^2)$ computational complexity that is entirely independent of sample size $n$.

\subsection{Numerical Stability Considerations and Limitations}

When converting uncentered raw moments of very high order ($r > 20$) for empirical datasets with large location offsets, binomial expansions can be susceptible to floating-point catastrophic cancellation due to alternating sum signs. For ultra-high-order empirical transformations ($r \ge 30$), combining the binomial transformation with arbitrary-precision arithmetic libraries or log-space moment representations \citep{dinardo2008} provides numerical safeguards.

\section{Conclusion and Future Work}

In this paper, we have presented a generalized binomial framework for transforming statistical moments of any order $r \in \mathbb{N}$ across arbitrary origins $a \to k$. By introducing the origin shift parameter $b = a - k$ and the symbolic operator representation $\mu'_r(k) \equiv (a + b)^r$, we established a unified method that encompasses raw-to-raw, raw-to-central, and central-to-raw conversions under a single mathematical principle.

We validated the framework theoretically across discrete (Binomial) and continuous (Normal) probability distributions and presented an open-source R package implementation (\texttt{convmoment}). Computational benchmarks demonstrate that transforming pre-computed moments using our algorithm achieves a $>250\times$ speedup over raw data recomputation on large datasets ($n = 100,000$) while remaining numerically exact.

Future extensions of this work include generalizing the binomial shift framework to \textbf{multivariate product moments}, joint cross-moments, and matrix-valued spatial moment invariants in higher-dimensional computer vision applications.

\section{Supplementary Material}

The following supplementary materials are available online:

\begin{enumerate}
\item \textbf{R Package \texttt{convmoment}:} Complete source code, documentation, and vignette demonstrating all package functions. Available at \url{https://github.com/yourusername/convmoment}.
\item \textbf{Replication Code:} Full R script reproducing all numerical examples, distribution validations (Binomial, Normal), empirical validation, and computational benchmarks presented in the paper.
\item \textbf{Extended Benchmark Results:} Additional performance comparisons for higher moment orders ($r = 10, 20, 50$) and larger sample sizes ($n = 10^6$).
\item \textbf{Vignette:} A comprehensive tutorial vignette included with the package showing basic usage, theoretical distribution examples, and matrix formulation.
\end{enumerate}

\bibliography{bibliography.bib}

\end{document}